\section{Tổng quan: nghịch lý của chiều sâu}
\label{sec:tongquan}

Lý thuyết xấp xỉ phổ quát khẳng định một mạng nơ-ron đủ rộng có thể xấp xỉ mọi
hàm liên tục trên tập compact với sai số tuỳ ý. Bổ sung thêm chiều sâu chỉ làm
lớp hàm biểu diễn được \emph{giàu hơn}: một mạng $L+1$ tầng luôn có thể mô phỏng
mạng $L$ tầng bằng cách đặt tầng thừa thành ánh xạ đồng nhất. Vì vậy, theo lý
thuyết, mạng sâu hơn không bao giờ được phép tệ hơn.

Thực nghiệm những năm 1990--2000 lại cho kết quả ngược lại. Khi tăng số tầng của
một mạng kết nối đầy đủ dùng hàm kích hoạt sigmoid từ $5$ lên $20$ tầng,
\emph{sai số huấn luyện} --- chứ không phải sai số kiểm tra --- tăng lên. Đây là
điểm mấu chốt cần tách bạch ngay từ đầu:

\begin{itemize}[leftmargin=*, itemsep=2pt]
  \item Nếu sai số huấn luyện thấp nhưng sai số kiểm tra cao, ta gặp vấn đề
    \textbf{tổng quát hoá} (quá khớp). Cách chữa là chính quy hoá, thêm dữ liệu,
    giảm dung lượng mô hình.
  \item Nếu \textbf{sai số huấn luyện} tự nó không giảm được, mô hình thậm chí
    không học nổi dữ liệu nó đang nhìn thấy. Đây là vấn đề \textbf{tối ưu hoá}.
    Thêm dữ liệu hay chính quy hoá đều vô ích.
\end{itemize}

Hiện tượng triệt tiêu gradient thuộc loại thứ hai. Mô hình có đủ dung lượng biểu
diễn, hàm mục tiêu có tồn tại nghiệm tốt, nhưng thuật toán hạ gradient không
truyền được tín hiệu học tới các tầng đầu. Nói cách khác, \emph{không gian giả
thuyết chứa nghiệm nhưng quỹ đạo tối ưu không tới được nó}.

\begin{infobox}{Phát biểu vấn đề trong một câu}
Trong mạng nơ-ron sâu, gradient của hàm mất mát theo tham số của tầng thứ $l$
tỉ lệ với một \emph{tích} gồm $L-l$ thừa số ma trận; nếu độ lớn đặc trưng của
mỗi thừa số lệch khỏi $1$, tích đó suy giảm hoặc bùng nổ theo cấp số nhân, khiến
các tầng gần đầu vào hoặc không nhận được tín hiệu học, hoặc nhận tín hiệu vô
nghĩa về mặt số học.
\end{infobox}

\subsection{Hai chế độ suy biến}

Gọi $\boldsymbol{g}^{[l]} = \partial\mathcal{L}/\partial \mathbf{W}^{[l]}$ là
gradient tại tầng $l$. Hai chế độ hỏng hóc là đối xứng nhau:

\textbf{Triệt tiêu gradient (vanishing gradient).} $\|\boldsymbol{g}^{[l]}\|
\to 0$ theo cấp số nhân khi $l$ giảm về phía đầu vào. Biểu hiện thực tế: mất mát
giảm rất nhanh trong vài epoch đầu rồi đứng yên ở một mức cao; kiểm tra trọng số
cho thấy các tầng cuối thay đổi rõ rệt còn các tầng đầu gần như giữ nguyên giá
trị khởi tạo. Đây là trường hợp chính của tài liệu này, và là trường hợp
\emph{nguy hiểm hơn} vì nó im lặng: không có cảnh báo, không có \texttt{NaN},
chỉ có một mô hình học kém mà ta dễ quy nhầm cho thiếu dữ liệu.

\textbf{Bùng nổ gradient (exploding gradient).} $\|\boldsymbol{g}^{[l]}\|
\to \infty$ theo cấp số nhân. Biểu hiện: mất mát nhảy vọt rồi trở thành
\texttt{NaN} hoặc \texttt{Inf} sau vài bước cập nhật. Chế độ này ồn ào nên dễ
phát hiện, và có một cách chữa cơ học trực tiếp là cắt ngưỡng gradient
\emph{(gradient clipping)}~\cite{Pascanu2013}.

Điểm sẽ được chứng minh ở Mục~\ref{sec:nguyennhan} là hai hiện tượng này
\textbf{không phải hai lỗi khác nhau}. Chúng là cùng một cơ chế --- tích ma trận
lặp --- quan sát ở hai phía của ngưỡng $1$.

\subsection{Cái bẫy mạng nông}

Hệ quả lịch sử của hiện tượng này là điều có thể gọi là \emph{cái bẫy mạng nông}.
Khi mạng sâu huấn luyện kém hơn mạng nông, kết luận tự nhiên --- và sai --- là
``chiều sâu không giúp ích cho bài toán này''. Người thực hành bèn quay lại kiến
trúc $2$--$3$ tầng, thấy nó hoạt động ổn định hơn, và củng cố niềm tin sai lầm
đó. Suốt gần hai thập kỷ, cộng đồng đã đọc một \emph{thất bại của thuật toán tối
ưu} thành một \emph{giới hạn của kiến trúc}.

Chẩn đoán đúng đến từ luận văn của Hochreiter~\cite{Hochreiter1991} và phân tích
hình thức của Bengio, Simard và Frasconi~\cite{BengioSimardFrasconi1994}: vấn đề
nằm ở cấu trúc nhân của lan truyền ngược, không nằm ở chiều sâu. Khi nguyên nhân
được hiểu đúng, các biện pháp khắc phục --- khởi tạo có kiểm soát phương
sai~\cite{GlorotBengio2010,He2015}, hàm kích hoạt không bão
hoà~\cite{He2015}, chuẩn hoá theo lô~\cite{IoffeSzegedy2015} và kết nối
tắt~\cite{He2016ResNet} --- đã cho phép huấn luyện mạng hàng trăm tầng.

\begin{notebox}
Trước khi kết luận ``mô hình của tôi cần thêm dữ liệu'', hãy kiểm tra sai số
\emph{huấn luyện}. Nếu chính nó không giảm, dữ liệu không phải là biến số cần
điều chỉnh.
\end{notebox}

\subsection{Cấu trúc tài liệu}

Mục~\ref{sec:cosotoan} thiết lập ký hiệu và suy diễn đầy đủ công thức lan truyền
ngược dạng ma trận. Mục~\ref{sec:nguyennhan} tách vấn đề thành hai nguyên nhân độc
lập --- độ bão hoà của hàm kích hoạt và bán kính phổ của ma trận trọng số --- rồi
hợp nhất chúng thành một đại lượng điều khiển duy nhất. Mục~\ref{sec:vidu} kiểm
chứng bằng một ví dụ số học ba tầng tính tay. Mục~\ref{sec:bptt} mở rộng sang mạng
hồi quy, nơi vấn đề nghiêm trọng hơn hẳn. Mục~\ref{sec:thucnghiem} đo đạc thực
nghiệm trên mạng $10$ tầng. Mục~\ref{sec:tongket} tổng kết các nguyên lý khắc phục.
